\chapter{Équations et mise en équation}
\label{corriges:06}
\section*{Exercice 6.1 -- Premier degré et cas particuliers}
La première équation donne $10x-5=3x+16$, puis $7x=21$, donc $x=3$.
La deuxième, multipliée par $6$, donne $2x-3=5$, donc $x=4$.
La troisième donne $3=5$, impossible : aucune solution.
La quatrième est une identité : tout réel est solution.

\section*{Exercice 6.2 -- Produits et carrés}
Par produit nul : $x=3/2$ ou $x=-4$.
Pour $x^2=45$, $x=\pm\sqrt{45}=\pm3\sqrt5$.
Pour $x^2=-4$, aucune solution réelle.
Enfin $x^2-7x=x(x-7)=0$, donc $x=0$ ou $x=7$. Diviser directement par $x$ ferait perdre $0$.

\section*{Exercice 6.3 -- Un rectangle à trouver}
Si $x>0$ est la largeur, la longueur vaut $x+7$. L'équation $2x+2(x+7)=58$ donne $4x=44$, donc $x=11$. Les dimensions sont $11$ et $18\unite{cm}$, et l'aire vaut $198\unite{cm^2}$. Vérification : $2(11+18)=58$.

\section*{Exercice 6.4 -- Seuil de rentabilité}
Pour $n$ séances : $C(n)=24+4n$, $S(n)=7n$. L'égalité donne $24=3n$, donc $n=8$. Pour un avantage strict, $24+4n<7n$ équivaut à $n>8$ ; il faut donc au moins $9$ séances entières. À $8$ séances, les deux coûts valent $56$ euros.
